% ============================================================
% 07_conclusion.tex  (revised 2026-10)
% ============================================================

This paper revisited an adaptive ensemble for photovoltaic forecasting with the aim of making
its evaluation honest and reproducible, and extended it with archived weather forecasts and
calibrated prediction intervals. Hourly forecasts for lead times of 1 to 24~h were evaluated in
power on two contrasting sites: the four uncurtailed sites of the Yulara off-grid mini-grid
(hot desert, 768~kW) and the NIST Ground array (humid subtropical, 271~kW).

\paragraph{Key findings.}
\begin{itemize}
  \item The Yulara series used in earlier studies, the Desert Gardens site, is curtailed around
  noon; we estimate that half of its available energy (31.7--59.1\%) was not delivered. Benchmarks
  on such a series forecast a dispatch rule rather than the solar resource.

  \item Archived GFS forecasts, used only when published before the forecast issue time, are the
  main source of skill: they lower the MAE of AHSE by 13\% at Yulara (36.18 $\to$ 31.32~kW,
  significant for each of five seeds) and by 26\% at NIST (27.55 $\to$ 20.39~kW), with gains
  growing with the lead time.

  \item The per-horizon-block selection with a day-level bootstrap guard combines models only
  where the validation evidence supports it. It beats the best single model of every seed by
  7--10\% at Yulara without weather forecasts and, when one member dominates (GFS-fed XGBoost),
  falls back to it; it was never significantly worse than the best member in hindsight.

  \item Out-of-fold, Mondrian and adaptive conformal calibration give 90\% intervals with
  0.90--0.92 coverage on the test periods; weather forecasts narrow them by 17\% at Yulara, and
  calibrating them by the GFS clearness index gives the best interval score at both sites.

  \item Evaluation choices matter: the MAE of the clearness index and the MAE alone can reverse
  the ranking of close models, so results should be reported in power with a squared-error skill
  score and day-level significance tests.
\end{itemize}

\paragraph{Future work.}
The work opens three directions linked to resilient microgrids. First, forecasting the
\emph{available} power of curtailed sites, reconstructed from uncurtailed neighbours, would give
the operator of an off-grid mini-grid the information needed to schedule its gas generators and
storage instead of reproducing past curtailment decisions. Second, ensemble weather predictions
and probabilistic members would allow the prediction intervals to reflect weather uncertainty
directly, for reserve sizing and risk-aware dispatch. Third, the honest selection protocol can be
coupled with physics-informed models of the PV array and the microgrid, to keep the forecasts
consistent with the physical limits that a dispatch controller has to respect.
